%% 2-13 公式に数を入れてみよう：代入　組5（第11回で書き足し）
%% 代入と式の値（×をもどす・累乗が先）、単位ごと代入する（換算は単位つきの量 60 s/min をかける。X11）、
%% マイナス・分数はかっこをつけて代入、℃と°Fの式の作り方、式を「かっこごと入れる」（第5巻 5-2・5-3「かっこごと入れるのが代入です（2-13）」）、
%% 代入をくり返して、はさんで調べる。ほしい量をとり出すのは 2-14。
\makeatletter
% 立てたものさし（下が0）  \bkiinruler{(x,y)}{1cmの長さmm}{目もりの最大cm}
\newcommand\bkiinruler[3]{\begin{scope}[shift={#1}]
  \fill[bkYellowBg, draw=bkYellow!80!black, line width=0.8pt, rounded corners=0.8mm] (0,-1.5) rectangle (9,#3*#2+1.5);
  \foreach \bkiin@k in {0,...,#3} {\draw[bkInk, line width=0.5pt] (0,\bkiin@k*#2) -- (2.4,\bkiin@k*#2);}
  \foreach \bkiin@k in {0,5,...,#3} {\draw[bkInk, line width=0.8pt] (0,\bkiin@k*#2) -- (3.6,\bkiin@k*#2);
    \node[font=\sffamily\fontsize{7}{8}\selectfont, anchor=west, inner sep=0.6pt] at (4,\bkiin@k*#2) {\bkiin@k};}
\end{scope}}
% つまむ指（親指と人さし指）  \bkiinfingers{(x,y)}{色}{線の形}
\newcommand\bkiinfingers[3]{\begin{scope}[shift={#1}]
  \filldraw[fill=#2, draw=bkSkin!60!black, line width=0.6pt, #3, rounded corners=1.6mm] (-6.5,-1.5) rectangle (0,1.9);
  \filldraw[fill=#2, draw=bkSkin!60!black, line width=0.6pt, #3, rounded corners=1.6mm] (9,-1.5) rectangle (15.5,1.9);
\end{scope}}
% 式を、まん中に1行で置く  \bkiinline{…}
\newcommand\bkiinline[1]{\par\addvspace{3pt}{\centering #1\par}\addvspace{3pt}}
% 数を入れる箱  \bkiinbox{(x,y)}{名前}{中の文字}{色}
\newcommand\bkiinbox[4]{\node[draw=#4, line width=1.1pt, fill=#4!8, rounded corners=0.8mm, minimum width=10mm, minimum height=9mm, font=\large, inner sep=1pt] (#2) at #1 {#3};}
% 数のカード  \bkiincard{(x,y)}{名前}{数}
\newcommand\bkiincard[3]{\node[draw=bkBlue, line width=0.9pt, fill=white, rounded corners=1mm, minimum height=7mm, minimum width=10mm, font=\normalsize, inner sep=1.5pt] (#2) at #1 {#3};}
\makeatother

\bkchapter{2}{13}{\R{公式}{こう|しき}に\R{数}{かず}を\R{入}{い}れてみよう：\R{代入}{だい|にゅう}}{\bklevel{2}}{\R{公式}{こう|しき}に\R{数}{かず}を\R{入}{い}れてみよう：\R{代入}{だい|にゅう}}
\bkmichishirube{2-2,2-10,2-12}{2-14,2-16,3-3,5-2,5-3}

\begin{goal}
\begin{itemize}[label={\color{bkTheme}●}, leftmargin=1.6\zw, itemsep=0pt, topsep=1pt]
\item \bklevel{1}\enspace \R{公式}{こう|しき}の\R{文字}{も|じ}に\R{数}{かず}を\R{代入}{だい|にゅう}して、\R{式}{しき}の\R{値}{あたい}を\R{求}{もと}められる。
\item \bklevel{2}\enspace \R{単位}{たん|い}ごと\R{代入}{だい|にゅう}し、\R{単位}{たん|い}をそろえて\R{計算}{けい|さん}できる。
\item \bklevel{2}\enspace マイナスの\R{数}{かず}・\R{分数}{ぶん|すう}・\R{式}{しき}を、かっこごと\R{代入}{だい|にゅう}できる。
\end{itemize}
\end{goal}

\begin{hatena}
\noindent\begin{minipage}[c]{0.65\linewidth}
\R{友}{とも}だちが\R{急}{きゅう}に\R{手}{て}をはなしたものさしを、すばやくつかみます。ものさしは、つかむまでに\qty{20}{cm}\R{落}{お}ちていました。\R{手}{て}をはなしてからつかむまでに、\R{何秒}{なん|びょう}かかったのかな？\par\vspace{4pt}
{\small \R{物}{もの}が\R{落}{お}ちはじめてから $t$ \R{秒間}{びょう|かん}に\R{落}{お}ちる\R{距離}{きょ|り}を $x$ m とすると、$x=5t^2$ と\R{表}{あらわ}せます。2-0の「1\R{秒}{びょう}ごとに\R{約}{やく}\qty{10}{m/s}ずつ\R{速}{はや}くなる」から\R{作}{つく}れる\R{式}{しき}です。なぜ5なのかは、この\R{章}{しょう}の「もっと\R{知}{し}りたい」で\R{確}{たし}かめます。}
\end{minipage}\hfill
\begin{minipage}[c]{0.32\linewidth}\centering
\begin{tikzpicture}[x=1mm, y=1mm]
  \bkiinruler{(0,0)}{1.3}{30}
  \fill[bkSkin, rounded corners=2mm] (-2,40.5) rectangle (11,46);
  \bkiinfingers{(0,0)}{bkSkin!50}{dashed}
  \bkiinfingers{(0,26)}{bkSkin}{solid}
  \node[font=\sffamily\bfseries\footnotesize, text=bkRed, anchor=west, align=left] at (16.2,26) {つかんだ\\\R{所}{ところ}};
  \node[font=\sffamily\footnotesize, text=bkGray, anchor=west] at (16.2,0) {はじめ};
  \node[font=\sffamily\footnotesize, text=bkGray, anchor=west] at (11.5,43.2) {はなす};
\end{tikzpicture}
\end{minipage}
\end{hatena}

\R{公式}{こう|しき}の\R{文字}{も|じ}に\R{数}{かず}を\R{入}{い}れると、\R{答}{こた}えが\R{出}{で}てきます。この\R{章}{しょう}では、\R{公式}{こう|しき}に\R{数}{かず}を\R{入}{い}れる\R{計算}{けい|さん}と、そのときに\R{気}{き}をつけることを\R{学}{まな}びます。\R{気}{き}をつけることは4つあります。×をもどすこと、\R{単位}{たん|い}をそろえること、かっこをつけること、\R{計算}{けい|さん}の\R{順番}{じゅん|ばん}を\R{守}{まも}ることです。

\clearpage
%% ---------------- やってみよう・思い出そう ----------------
\begin{yattemiyou}[ものさしキャッチで、\R{反応時間}{はん|のう|じ|かん}をはかろう]
\begin{enumerate}[label={\bkcircled[bkGreen]{\arabic*}}, leftmargin=2.4\zw, itemsep=1pt, topsep=2pt]
\item 30 cmのものさしを、0の\R{目}{め}もりを\R{下}{した}にして、\R{友}{とも}だちに\R{持}{も}ってもらう。
\item 0の\R{所}{ところ}に\R{指}{ゆび}をそえて\R{待}{ま}ち、\R{友}{とも}だちが\R{合図}{あい|ず}なしで\R{手}{て}をはなしたら、つかむ。
\item つかんだ\R{所}{ところ}の\R{目}{め}もり（cm）を\R{読}{よ}む。5\R{回}{かい}やって、\R{表}{ひょう}に\R{書}{か}こう。
\item \R{何}{なん}cmなら\R{何秒}{なん|びょう}か、\R{数学}{すう|がく}の\R{道具}{どう|ぐ}①の\R{表}{ひょう}で\R{調}{しら}べよう。
\end{enumerate}
\begin{bkcenter}
\begin{bktable}{colspec={Q[l,m,wd=34mm] *{5}{Q[c,m,wd=13mm]}}, rowsep=1.5pt, colsep=4pt, cells={font=\small}, row{1}={bg=bkGreen!14}}
 & 1\R{回目}{かい|め} & 2\R{回目}{かい|め} & 3\R{回目}{かい|め} & 4\R{回目}{かい|め} & 5\R{回目}{かい|め} \\
つかんだ\R{所}{ところ}（cm） & & & & & \\
\R{反応時間}{はん|のう|じ|かん}（\R{秒}{びょう}） & & & & & \\
\end{bktable}
\end{bkcenter}
{\small ものさしが\R{落}{お}ちるとき、まわりの\R{人}{ひと}にあたらないように\R{気}{き}をつけよう。}
\end{yattemiyou}

\begin{omoidasou}
\bkrecall{2-12}{$5t^2$ は $5\times t\times t$。×をはぶいて\R{書}{か}いている。}
\bkrecall{2-11}{\R{計算}{けい|さん}の\R{順番}{じゅん|ばん}：かっこの\R{中}{なか}が\R{先}{さき}。×÷は＋−より\R{先}{さき}。}
\bkrecall{2-2}{$(-3)^2=(-3)\times(-3)=9$。マイナスの\R{数}{かず}を2\R{回}{かい}かけるとプラス。}
\bkrecall{2-10}{\R{換算}{かん|さん}は「1をかける」。\qty{60}{s/min}は「1\R{分}{ぷん}あたり60\R{秒}{びょう}」。}
\bkrecall{2-11}{\R{文字}{も|じ}は、どんな\R{数}{かず}でも\R{入}{はい}る\R{箱}{はこ}。}
\end{omoidasou}

\bkfukidashi{sora}{つかむのが\R{速}{はや}い\R{人}{ひと}ほど、ものさしの\R{小}{ちい}さい\R{目}{め}もりの\R{所}{ところ}でつかめるね。\R{何}{なん}cmなら\R{何秒}{なん|びょう}なのか、\R{式}{しき}で\R{計算}{けい|さん}してみよう。}

\begin{column}{\R{反応時間}{はん|のう|じ|かん}って、どれくらい？}
\noindent\begin{minipage}[c]{0.56\linewidth}
\R{目}{め}で「\R{落}{お}ちた！」と\R{見}{み}てから、\R{指}{ゆび}が\R{動}{うご}くまでには、\R{少}{すこ}し\R{時間}{じ|かん}がかかります。\R{目}{め}からの\R{知}{し}らせが、\R{神経}{しん|けい}を\R{通}{とお}って\R{脳}{のう}にとどき、\R{脳}{のう}から\R{指}{ゆび}の\R{筋肉}{きん|にく}に\R{命令}{めい|れい}が\R{行}{い}くからです。この\R{時間}{じ|かん}を\textbf{\R{反応時間}{はん|のう|じ|かん}}といいます。ふつうは、およそ0.15\R{秒}{びょう}から0.3\R{秒}{びょう}です。\R{自動車}{じ|どう|しゃ}の\R{運転}{うん|てん}では、ブレーキをふむまでのこの\R{時間}{じ|かん}にも\R{車}{くるま}が\R{進}{すす}むので、\R{車}{くるま}と\R{車}{くるま}のあいだをあけて\R{走}{はし}ります。
\end{minipage}\hfill
\begin{minipage}[c]{0.42\linewidth}\centering
\begin{tikzpicture}[x=1mm, y=1mm]
  \node[bk node, minimum width=11mm] (e) at (0,0) {\R{目}{め}};
  \node[bk node, minimum width=11mm] (b) at (19,0) {\R{脳}{のう}};
  \node[bk node, minimum width=11mm] (f) at (38,0) {\R{指}{ゆび}};
  \draw[bk link] (e) -- (b);
  \draw[bk link] (b) -- (f);
  \node[font=\sffamily\footnotesize, anchor=north] at (19,-7) {\R{合計}{ごう|けい}で\R{約}{やく}0.2\R{秒}{びょう}};
\end{tikzpicture}
\end{minipage}
\end{column}

%% ---------------- 数学の道具 ----------------
\bksection{doguu}{\R{数学}{すう|がく}の\R{道具}{どう|ぐ}}

\bksubsection{1}{\R{文字}{も|じ}に\R{数}{かず}を\R{入}{い}れる：\R{代入}{だい|にゅう}}

\leavevmode\label{bk:2-13:dainyu}\bkhatenaref のものさしは、\R{何秒}{なん|びょう}で\qty{20}{cm}\R{落}{お}ちるのでしょう。まず、0.2\R{秒}{びょう}で\R{何}{なん}cm\R{落}{お}ちるかを\R{計算}{けい|さん}してみます。$x=5t^2$ の $t$ に、0.2を\R{入}{い}れればよいですね。

2-11で、\R{文字}{も|じ}は「どんな\R{数}{かず}でも\R{入}{はい}る\R{箱}{はこ}」でした。2-12で、$5t^2$ は $5\times t\times t$ を\R{短}{みじか}く\R{書}{か}いたものでした。そこで、はぶいていた×をもとにもどして、$t$ の\R{箱}{はこ}を2つならべます。それぞれの\R{箱}{はこ}に、0.2を\R{入}{い}れます。

\begin{bkfigure}
\begin{tikzpicture}[x=1mm, y=1mm]
  \node[font=\large, anchor=east] at (21,0) {$x=5\times$};
  \bkiinbox{(28,0)}{b1}{$t$}{bkTheme}
  \node[font=\large] at (38,0) {$\times$};
  \bkiinbox{(48,0)}{b2}{$t$}{bkTheme}
  \bkiincard{(28,17)}{c1}{0.2}
  \bkiincard{(48,17)}{c2}{0.2}
  \draw[bk link] (c1) -- (b1);
  \draw[bk link] (c2) -- (b2);
  \draw[bk arrow, bkGray] (60,0) -- (72,0);
  \node[font=\large, anchor=west] at (75,0) {$x=5\times0.2\times0.2$};
  \node[font=\sffamily\footnotesize, anchor=south west] at (56,17) {\R{同}{おな}じ\R{文字}{も|じ}には、\R{同}{おな}じ\R{数}{かず}を\R{入}{い}れる};
\end{tikzpicture}
\bkfigcaption{×をもどして、$t$ の\R{箱}{はこ}に0.2を\R{入}{い}れる。}
\end{bkfigure}

\R{同}{おな}じ\R{文字}{も|じ} $t$ が2\R{回}{かい}\R{出}{で}てくるときは、どちらにも\R{同}{おな}じ0.2を\R{入}{い}れます。\R{同}{おな}じ\R{式}{しき}の\R{中}{なか}の\R{同}{おな}じ\R{文字}{も|じ}は、\R{同}{おな}じ\R{量}{りょう}（ここでは\R{同}{おな}じ\R{時間}{じ|かん}）を\R{表}{あらわ}しているからです。

\R{次}{つぎ}に、\R{計算}{けい|さん}します。0.2×0.2は、2×2＝4で、\R{小数点}{しょう|すう|てん}の\R{下}{した}のけたが\R{合}{あ}わせて2つなので0.04です（$\dfrac{2}{10}\times\dfrac{2}{10}=\dfrac{4}{100}$ と\R{考}{かんが}えてもよい。2-4）。
\bkiinline{$x=5\times0.2\times0.2=5\times0.04=0.2$}
\R{答}{こた}えの\R{単位}{たん|い}はmなので、\qty{0.2}{m}です。cmに\R{直}{なお}すには、「1 mあたり\qty{100}{cm}」を\R{表}{あらわ}す \qty{100}{cm/m} をかけます（2-10）。mどうしが\R{約分}{やく|ぶん}されて\R{消}{き}え、cmが\R{残}{のこ}ります。
\bkiinline{$\qty{0.2}{m}\times\qty{100}{cm/m}=0.2\,\cancel{\mathrm{m}}\times\dfrac{100\,\mathrm{cm}}{1\,\cancel{\mathrm{m}}}=\qty{20}{cm}$}
つまり、0.2\R{秒}{びょう}でちょうど\qty{20}{cm}\R{落}{お}ちます。

このように、\R{式}{しき}の\R{文字}{も|じ}に\R{数}{かず}を\R{入}{い}れることを\textbf{\R{代入}{だい|にゅう}}といい、\bkidx{だいにゅう}{\R{代入}{だい|にゅう}}\R{代入}{だい|にゅう}して\R{計算}{けい|さん}した\R{答}{こた}えを\textbf{\R{式}{しき}の\R{値}{あたい}}といいます。\bkidx{しきのあたい}{\R{式}{しき}の\R{値}{あたい}}\R{代入}{だい|にゅう}の\R{手順}{て|じゅん}をまとめると、\R{次}{つぎ}のようになります。

\begin{bkcenter}
\begin{tikzpicture}[x=1mm, y=1mm]
  \node[bk node, minimum width=34mm, minimum height=13mm] (a) at (0,0) {① はぶいていた\\×をもどす};
  \node[bk node, minimum width=34mm, minimum height=13mm] (b) at (44,0) {② \R{文字}{も|じ}の\R{所}{ところ}に\\\R{数}{かず}を\R{書}{か}く};
  \node[bk node, minimum width=34mm, minimum height=13mm] (c) at (88,0) {③ \R{計算}{けい|さん}の\R{順番}{じゅん|ばん}を\\\R{守}{まも}って\R{計算}{けい|さん}する};
  \draw[bk link] (a) -- (b); \draw[bk link] (b) -- (c);
\end{tikzpicture}
\end{bkcenter}

\textbf{なぜ×をもどすの？}\quad ×をもどさずに、$5t^2$ の $t$ のところにそのまま0.2を\R{書}{か}くと、「$50.2^2$」になってしまいます。これでは、「50.2の2\R{乗}{じょう}」と\R{読}{よ}めてしまいますね。$5t$ に $t=3$ を\R{入}{い}れて「53」と\R{書}{か}くと、\R{五十三}{ご|じゅう|さん}になってしまいます。\R{文字}{も|じ}のときははぶいてよかった×も、\R{数}{かず}を\R{入}{い}れたら\R{必}{かなら}ずもどします。

\textbf{なぜ2\R{乗}{じょう}を\R{先}{さき}に\R{計算}{けい|さん}するの？}\quad $5\times0.2^2$ の「2」は、すぐ\R{前}{まえ}の0.2だけについています。$5t^2$ は $5\times t\times t$ で、2\R{回}{かい}かけるのは $t$ だけだからです。だから、2\R{乗}{じょう}は「$t\times t$ のひとまとまり」として、×÷より\R{先}{さき}に\R{計算}{けい|さん}します。もし $5\times0.2$ を\R{先}{さき}にして $(5\times0.2)^2=1^2=1$ とすると、\qty{1}{m}、つまり\qty{100}{cm}\R{落}{お}ちたことになります。\qty{30}{cm}のものさしでは、ありえませんね。

\begin{bkcenter}
{\setlength\fboxsep{4pt}\fcolorbox{bkTheme}{white}{\R{計算}{けい|さん}の\R{順番}{じゅん|ばん}：かっこの\R{中}{なか}\enspace →\enspace 2\R{乗}{じょう}などの\R{累乗}{るい|じょう}\enspace →\enspace ×÷\enspace →\enspace ＋−}}
\end{bkcenter}

\R{同}{おな}じように、$t$ にいろいろな\R{数}{かず}を\R{代入}{だい|にゅう}して、\R{表}{ひょう}を\R{作}{つく}りましょう。\R{時間}{じ|かん}を0.02\R{秒}{びょう}ずつふやしています。\R{電卓}{でん|たく}を\R{使}{つか}ってもかまいません。

\begin{bkcenter}
\begin{bktable}{colspec={Q[l,m,wd=27mm] *{8}{Q[c,m,wd=9mm]}}, rowsep=1pt, colsep=3pt, cells={font=\small}}
\R{時間}{じ|かん} $t$（\R{秒}{びょう}） & 0.10 & 0.12 & 0.14 & 0.16 & 0.18 & 0.20 & 0.22 & 0.24 \\
$5t^2$（m） & 0.05 & 0.072 & 0.098 & 0.128 & 0.162 & 0.2 & 0.242 & 0.288 \\
\R{落}{お}ちた\R{長}{なが}さ（cm） & 5 & 7.2 & 9.8 & 12.8 & 16.2 & \SetCell{bg=bkYellow!40}20 & 24.2 & 28.8 \\
\end{bktable}
\bktabcaption{$x=5t^2$ に $t$ を\R{代入}{だい|にゅう}して\R{作}{つく}った\R{表}{ひょう}（cmの\R{行}{ぎょう}は、mの\R{行}{ぎょう}に \qty{100}{cm/m} をかけた）}
\end{bkcenter}

\bkhatenaref の\R{答}{こた}えは、\qty{20}{cm}\R{落}{お}ちたから\textbf{0.2\R{秒}{びょう}}です。\R{表}{ひょう}を\R{見}{み}ると、\R{時間}{じ|かん}が0.10\R{秒}{びょう}から0.20\R{秒}{びょう}へ2\R{倍}{ばい}になると、\R{落}{お}ちた\R{長}{なが}さは\qty{5}{cm}から\qty{20}{cm}へ4\R{倍}{ばい}になっています。2\R{乗}{じょう}があるので、$2\times2=4$\R{倍}{ばい}になるのです（2-15でくわしく）。

\R{表}{ひょう}の\R{数}{かず}をものさしに\R{書}{か}きこむと、\R{反応時間}{はん|のう|じ|かん}をすぐに\R{読}{よ}める「\R{反応時間}{はん|のう|じ|かん}ものさし」ができます。

\begin{bkfigure}
\begin{tikzpicture}[x=1mm, y=1mm]
  \fill[bkYellowBg, draw=bkYellow!80!black, line width=0.8pt, rounded corners=1mm] (-3,-1) rectangle (123,13);
  \foreach \k in {0,...,30} {\draw[bkInk, line width=0.5pt] (\k*4,13) -- (\k*4,10.6);}
  \foreach \k in {0,5,...,30} {\draw[bkInk, line width=0.8pt] (\k*4,13) -- (\k*4,9.4); \node[font=\sffamily\footnotesize, anchor=north] at (\k*4,8.8) {\k};}
  \node[font=\sffamily\footnotesize, anchor=west] at (124,10.5) {cm};
  \foreach \t/\c in {0.10/5, 0.14/9.8, 0.18/16.2, 0.22/24.2} {\draw[bkRed, line width=0.9pt] (\c*4,-1) -- (\c*4,2.6); \node[font=\sffamily\bfseries\footnotesize, text=bkRed, anchor=north] at (\c*4,-1.8) {\t};}
  \foreach \t/\c in {0.12/7.2, 0.16/12.8, 0.20/20, 0.24/28.8} {\draw[bkRed, line width=0.9pt] (\c*4,-1) -- (\c*4,2.6); \node[font=\sffamily\bfseries\footnotesize, text=bkRed, anchor=north] at (\c*4,-6.3) {\t};}
  \node[font=\sffamily\footnotesize, text=bkRed, anchor=west] at (124,-4) {\R{秒}{びょう}};
\end{tikzpicture}
\bkfigcaption{\R{反応時間}{はん|のう|じ|かん}ものさし。\R{上}{うえ}の\R{黒}{くろ}い\R{数}{かず}はcm、\R{下}{した}の\R{赤}{あか}い\R{数}{かず}は\R{秒}{びょう}。\R{赤}{あか}い\R{線}{せん}の\R{間隔}{かん|かく}は、\R{下}{した}（0の\R{方}{ほう}）ほどせまい。}
\end{bkfigure}

\R{赤}{あか}い\R{線}{せん}は、\R{同}{おな}じ0.02\R{秒}{びょう}ずつなのに、\R{右}{みぎ}へ\R{行}{い}くほど\R{間隔}{かん|かく}が\R{広}{ひろ}くなっています。\R{落}{お}ちるものさしは、だんだん\R{速}{はや}くなるからです。

\Needspace{16\baselineskip}\begin{motto}[2]
\textbf{なぜ $x=5t^2$ なの？}\par
\noindent\begin{minipage}[c]{0.52\linewidth}
\R{落}{お}ちはじめの\R{速}{はや}さは0です。1\R{秒}{びょう}ごとに\R{約}{やく}\qty{10}{m/s}ずつ\R{速}{はや}くなるので（2-0）、$t$ \R{秒後}{びょう|ご}の\R{速}{はや}さは $10\times t$（m/s）です。\R{速}{はや}さは\R{同}{おな}じ\R{割合}{わり|あい}でふえるので、0\R{秒}{びょう}から $t$ \R{秒}{びょう}までの\R{平均}{へい|きん}の\R{速}{はや}さ（2-6）は、はじめの0とおわりの $10t$ の、ちょうどまん\R{中}{なか}です。
\[ (0+10t)\div2=5t\ \ (\mathrm{m/s}) \]
\R{道}{みち}のり＝\R{速}{はや}さ×\R{時間}{じ|かん}（2-3）なので、$x=5t\times t=5t^2$。5は、10の\R{半分}{はん|ぶん}だったのです。
\end{minipage}\hfill
\begin{minipage}[c]{0.45\linewidth}\centering
\begin{tikzpicture}
\begin{axis}[bkplot, width=42mm, height=46mm, xmin=0, xmax=2.5, ymin=0, ymax=24, xtick={0,1,2}, ytick={0,10,20},
  xlabel={\R{時間}{じ|かん}（\R{秒}{びょう}）}, ylabel={\R{速}{はや}さ（m/s）}]
  \addplot[fill=bkTheme!25, draw=none] coordinates {(0,0) (2,20) (2,0) (0,0)};
  \addplot[bkcurve, domain=0:2.3] {10*x};
  \addplot[bkBlue, dashed, line width=0.9pt] coordinates {(0,10) (2,10) (2,0)};
  \node[font=\sffamily\scriptsize, anchor=south west, inner sep=1pt] at (axis cs:0.05,11.8) {\R{平均}{へい|きん}10};
\end{axis}
\end{tikzpicture}
\end{minipage}
\par\vspace{3pt}
\R{右}{みぎ}のグラフは $t=2$ のときです。\R{色}{いろ}の\R{三角}{さん|かく}の\R{面積}{めん|せき}が\R{落}{お}ちた\R{距離}{きょ|り}を\R{表}{あらわ}します（\R{速}{はや}さが\R{変}{か}わるときもグラフの\R{面積}{めん|せき}が\R{道}{みち}のりになることは、5-9で\R{確}{たし}かめます）。\R{三角}{さん|かく}の\R{面積}{めん|せき}は、\R{高}{たか}さを\R{半分}{はん|ぶん}の10にした\R{長方形}{ちょう|ほう|けい}（\R{点線}{てん|せん}）と\R{同}{おな}じで、$10\times2=20$（m）。$5\times2^2=20$ と\R{合}{あ}います。ほんとうは1\R{秒}{びょう}ごとに\qty{9.8}{m/s}ずつ\R{速}{はや}くなるので、くわしい\R{式}{しき}はその\R{半分}{はん|ぶん}の $4.9t^2$ です（5-5）。
\end{motto}

\bksubsection{2}{\R{単位}{たん|い}ごと\R{代入}{だい|にゅう}する}

\leavevmode\label{bk:2-13:tani}$x=vt$（\R{道}{みち}のり＝\R{速}{はや}さ×\R{時間}{じ|かん}）に、$v=\qty{5}{m/s}$、$t=\qty{4}{s}$ を\R{代入}{だい|にゅう}します。\R{数}{かず}だけでなく、\textbf{\R{単位}{たん|い}もいっしょに}\R{代入}{だい|にゅう}してみましょう。
\bkiinline{$x=\qty{5}{m/s}\times\qty{4}{s}=(5\times4)\,\dfrac{\mathrm{m}}{\cancel{\mathrm{s}}}\times\cancel{\mathrm{s}}=\qty{20}{m}$}
\R{数}{かず}は\R{数}{かず}どうし（$5\times4=20$）、\R{単位}{たん|い}は\R{単位}{たん|い}どうしで\R{計算}{けい|さん}します（2-10）。m/sは $\dfrac{\mathrm{m}}{\mathrm{s}}$ なので、sをかけると\R{分母}{ぶん|ぼ}のsと\R{約分}{やく|ぶん}されて\R{消}{き}え、mが\R{残}{のこ}ります。\R{道}{みち}のりの\R{単位}{たん|い}mになったので、\R{式}{しき}の\R{使}{つか}い\R{方}{かた}が\R{正}{ただ}しそうだとわかります。

\textbf{\R{単位}{たん|い}がそろっていないとき}\quad \R{時間}{じ|かん}が「2\R{分}{ふん}」のときは、どうでしょう。\R{数}{かず}だけを\R{入}{い}れて $5\times2=10$ とすると、まちがいです。\R{単位}{たん|い}をつけて\R{入}{い}れると、$\qty{5}{m/s}\times\qty{2}{min}$ となり、\R{秒}{びょう}（s）と\R{分}{ふん}（min）で\R{単位}{たん|い}がそろっていないことに\R{気}{き}づけます。sとminは\R{約分}{やく|ぶん}できないので、mだけが\R{残}{のこ}りません。

\noindent\begin{minipage}[c]{0.52\linewidth}
そこで、\R{先}{さき}に2\R{分}{ふん}を\R{秒}{びょう}に\R{換算}{かん|さん}します。2-10で\R{学}{まな}んだように、1\R{分}{ぷん}の\R{中}{なか}には60\R{秒}{びょう}があります。「1\R{分}{ぷん}あたり60\R{秒}{びょう}」は\R{単位}{たん|い}つきの\R{量}{りょう}で、\qty{60}{s/min}と\R{書}{か}きます（minはminute、sはsecondの\R{頭文字}{かしら|も|じ}）。\qty{60}{s}と\qty{1}{min}は\R{同}{おな}じ\R{長}{なが}さなので、$\dfrac{60\,\mathrm{s}}{1\,\mathrm{min}}$ は1と\R{同}{おな}じです。1をかけても、\R{量}{りょう}は\R{変}{か}わりません。
\end{minipage}\hfill
\begin{minipage}[c]{0.45\linewidth}\centering
\begin{tikzpicture}[x=1mm, y=1mm]
  \fill[bkTheme!15] (0,6) rectangle (48,12);
  \draw[bkInk, line width=0.8pt] (0,6) rectangle (48,12);
  \node[font=\sffamily\footnotesize] at (24,9) {\qty{1}{min}};
  \foreach \k in {0,...,60} {\draw[bkInk, line width=0.3pt] ({\k*0.8},0) -- ({\k*0.8},{ifthenelse(mod(\k,10)==0,3,1.6)});}
  \draw[bkInk, line width=0.8pt] (0,0) -- (48,0);
  \foreach \k in {0,10,...,60} \node[font=\sffamily\scriptsize, anchor=north] at ({\k*0.8},-0.8) {\k};
  \node[font=\sffamily\scriptsize, anchor=west] at (49,-2.2) {s};
\end{tikzpicture}
\bkfigcaption{1\R{分}{ぷん}の\R{中}{なか}に、60\R{秒}{びょう}がならぶ。}
\end{minipage}
\bkiinline{$\qty{2}{min}\times\qty{60}{s/min}=2\,\cancel{\mathrm{min}}\times\dfrac{60\,\mathrm{s}}{1\,\cancel{\mathrm{min}}}=\qty{120}{s}$}
minどうしが\R{約分}{やく|ぶん}されて\R{消}{き}え、sが\R{残}{のこ}りました。これで\R{単位}{たん|い}がそろったので、\R{代入}{だい|にゅう}します。
\bkiinline{$x=\qty{5}{m/s}\times\qty{120}{s}=\qty{600}{m}$}

\R{反対}{はん|たい}に、\R{秒}{びょう}を\R{分}{ふん}に\R{直}{なお}すときは、\qty{60}{s/min}でわります。\qty{120}{s}の\R{中}{なか}に、60\R{秒}{びょう}のまとまり（1\R{分}{ぷん}）がいくつあるかを\R{求}{もと}めるのです。
\bkiinline{$\qty{120}{s}\div\qty{60}{s/min}=120\,\cancel{\mathrm{s}}\times\dfrac{1\,\mathrm{min}}{60\,\cancel{\mathrm{s}}}=\qty{2}{min}$}
\R{時間}{じ|かん}・\R{分}{ふん}・\R{秒}{びょう}の\R{換算}{かん|さん}に\R{使}{つか}う\R{量}{りょう}は、\R{次}{つぎ}の3つです（hはhourの\R{頭文字}{かしら|も|じ}）。どれも「1あたり」の\R{量}{りょう}です。
\begin{bkcenter}
\begin{bktable}{colspec={Q[l,m,wd=30mm] Q[l,m,wd=44mm] Q[l,m,wd=42mm]}, rowsep=1pt, cells={font=\small}, row{1}={font=\sffamily\bfseries\small}}
\R{単位}{たん|い}つきの\R{量}{りょう} & \R{意味}{い|み} & なぜその\R{数}{かず}？ \\
\qty{60}{s/min} & 1\R{分}{ぷん}あたり60\R{秒}{びょう} & 1\R{分}{ぷん}は60\R{秒}{びょう}と\R{決}{き}めた \\
\qty{60}{min/h} & 1\R{時間}{じ|かん}あたり60\R{分}{ぷん} & 1\R{時間}{じ|かん}は60\R{分}{ぷん}と\R{決}{き}めた \\
\qty{3600}{s/h} & 1\R{時間}{じ|かん}あたり3600\R{秒}{びょう} & $60\times60=3600$ \\
\end{bktable}
\bktabcaption{\R{時間}{じ|かん}の\R{換算}{かん|さん}に\R{使}{つか}う\R{量}{りょう}（2-10）}
\end{bkcenter}
\R{長}{なが}さなら、\qty{100}{cm/m}（1 mあたり100 cm）や\qty{1000}{m/km}（1 kmあたり1000 m）を\R{使}{つか}います。

\textbf{\R{式}{しき}の\R{中}{なか}の\R{数}{かず}にも\R{単位}{たん|い}がある}\quad $x=5t^2$ の5にも、じつは\R{単位}{たん|い}があります。$x$ はm、$t^2$ は $\mathrm{s}\times\mathrm{s}=\mathrm{s^2}$ なので、5の\R{単位}{たん|い}が\unit{m/s^2}（2-10）なら、\R{両側}{りょう|がわ}の\R{単位}{たん|い}が\R{合}{あ}います。\R{単位}{たん|い}ごと\R{代入}{だい|にゅう}して\R{確}{たし}かめましょう。
\bkiinline{$\qty{5}{m/s^2}\times(\qty{0.2}{s})^2=\qty{5}{m/s^2}\times\qty{0.04}{s^2}=\qty{0.2}{m}$}
\unit{s^2}が\R{約分}{やく|ぶん}されて\R{消}{き}えて、mが\R{残}{のこ}りました。

\bksubsection{3}{マイナスの\R{数}{かず}・\R{分数}{ぶん|すう}は、かっこをつけて\R{代入}{だい|にゅう}する}

\leavevmode\label{bk:2-13:kakko}$2x^2$ に $x=-3$ を\R{代入}{だい|にゅう}します。$2x^2$ は $2\times x\times x$ なので、$x$ の\R{箱}{はこ}2つに $-3$ をまるごと\R{入}{い}れます。$-3$ は「マイナス」と「3」に\R{分}{わ}けられない1つの\R{数}{かず}なので、かっこでひとまとまりにして\R{入}{い}れます。
\bkiinline{$2x^2=2\times(-3)^2=2\times(-3)\times(-3)=2\times9=18$}
$(-3)\times(-3)=9$ は、マイナスどうしのかけ\R{算}{ざん}でプラスになるからです（2-2）。

\textbf{かっこをつけないと？}\quad かっこをつけずに $2\times-3^2$ と\R{書}{か}くと、2つおかしなことがおきます。1つめは、×と−がならんで、ひき\R{算}{ざん}なのかマイナスの\R{数}{かず}なのかわからないことです。2つめは、2\R{乗}{じょう}の「2」が、すぐ\R{前}{まえ}の3だけにつくことです。\R{数学}{すう|がく}では、$-3^2$ は「$3^2$ にマイナスをつけたもの」、つまり $-(3\times3)=-9$ と\R{読}{よ}むやくそくです。これで\R{計算}{けい|さん}すると $2\times(-9)=-18$ となり、\R{正}{ただ}しい18とちがってしまいます。

\begin{bkcenter}
\begin{bktable}{colspec={Q[c,m,wd=20mm] Q[l,m,wd=44mm] Q[c,m,wd=12mm]}, rowsep=1.5pt, cells={font=\small}, row{1}={font=\sffamily\bfseries\small}}
\R{書}{か}き\R{方}{かた} & \R{意味}{い|み} & \R{値}{あたい} \\
$(-3)^2$ & $(-3)\times(-3)$ & 9 \\
$-3^2$ & $-(3\times3)$ & $-9$ \\
\end{bktable}
\bktabcaption{かっこがあるかないかで、\R{値}{あたい}がちがう。}
\end{bkcenter}

\noindent\begin{minipage}[c]{0.62\linewidth}
\textbf{\R{分数}{ぶん|すう}を2\R{乗}{じょう}するとき}\quad \R{分数}{ぶん|すう}も、かっこをつけて\R{代入}{だい|にゅう}します。$2x^2$ に $x=\dfrac{1}{2}$ を\R{代入}{だい|にゅう}すると、
\[ 2\times\left(\frac{1}{2}\right)^2=2\times\frac{1}{2}\times\frac{1}{2}=2\times\frac{1}{4}=\frac{1}{2} \]
$\dfrac{1}{2}\times\dfrac{1}{2}=\dfrac{1}{4}$ は、\R{右}{みぎ}の\R{図}{ず}で\R{確}{たし}かめられます。1\R{辺}{ぺん}1の\R{正方形}{せい|ほう|けい}のたてと\R{横}{よこ}を\R{半分}{はん|ぶん}にすると、4つに\R{分}{わ}けた1つ\R{分}{ぶん}になります（2-4）。かっこがないと、$\dfrac{1^2}{2}$ と\R{区別}{く|べつ}がつきません。
\end{minipage}\hfill
\begin{minipage}[c]{0.35\linewidth}\centering
\begin{tikzpicture}[x=1mm, y=1mm]
  \fill[bkTheme!35] (0,15) rectangle (15,30);
  \draw[bkInk, line width=0.9pt] (0,0) rectangle (30,30);
  \draw[bkInk, line width=0.5pt, dashed] (15,0) -- (15,30) (0,15) -- (30,15);
  \node[font=\small] at (7.5,22.5) {$\dfrac{1}{4}$};
  \node[font=\sffamily\footnotesize, anchor=south] at (7.5,30.8) {$\dfrac{1}{2}$};
  \node[font=\sffamily\footnotesize, anchor=east] at (-0.8,22.5) {$\dfrac{1}{2}$};
  \node[font=\sffamily\footnotesize, anchor=north] at (15,-0.8) {1};
\end{tikzpicture}
\end{minipage}
\par\vspace{6pt}
\textbf{\R{温度}{おん|ど}の\R{換算}{かん|さん}}\quad アメリカでは、\R{温度}{おん|ど}を°F（\textbf{\R{華氏}{か|し}}）\bkidx{かし}{\R{華氏}{か|し}}で\R{表}{あらわ}します。\R{日本}{に|ほん}で\R{使}{つか}う℃と°Fでは、\R{目}{め}もりの0の\R{位置}{い|ち}も、1\R{目}{め}もりの\R{大}{おお}きさもちがいます。\R{下}{した}の\R{図}{ず}のように、\R{水}{みず}がこおる\R{温度}{おん|ど}は\qty{0}{\degreeCelsius}で32°F、\R{水}{みず}がふっとうする\R{温度}{おん|ど}は\qty{100}{\degreeCelsius}で212°Fです。

\begin{bkfigure}
\begin{tikzpicture}[x=1mm, y=1mm]
  % ℃（上）と°F（下）の目もり。C=-10 が x=1、C=100 が x=100（1℃=0.9mm）
  \draw[bkInk, line width=0.9pt] (0,12) -- (104,12);
  \draw[bkInk, line width=0.9pt] (0,0) -- (104,0);
  \foreach \c in {-10,0,...,100} {\pgfmathsetmacro\xx{10+\c*0.9}
    \draw[bkInk, line width=0.6pt] (\xx,12) -- (\xx,14.5);
    \node[font=\sffamily\scriptsize, anchor=south] at (\xx,15) {$\c$};
    \draw[bkInk, line width=0.6pt] (\xx,0) -- (\xx,-2.5);
    \pgfmathtruncatemacro\ff{round(1.8*\c+32)}
    \node[font=\sffamily\scriptsize, anchor=north] at (\xx,-3) {\ff};}
  \node[font=\sffamily\footnotesize, anchor=west] at (106,13) {℃};
  \node[font=\sffamily\footnotesize, anchor=west] at (106,-1) {°F};
  \draw[bkBlue, line width=0.6pt, densely dashed] (10,2) -- (10,10);
  \draw[bkRed, line width=0.6pt, densely dashed] (100,2) -- (100,10);
  \node[font=\sffamily\footnotesize, text=bkBlue!80!black, fill=white, inner sep=1pt] at (17,6) {こおる};
  \node[font=\sffamily\footnotesize, text=bkRed!80!black, fill=white, inner sep=1pt] at (91,6) {ふっとう};
\end{tikzpicture}
\bkfigcaption{\R{上}{うえ}が℃、\R{下}{した}が°F。℃が10ふえるごとに、°Fは18ふえる。}
\end{bkfigure}

\begin{enumerate}[label={\bkcircled[bkTheme]{\arabic*}}, leftmargin=2.4\zw, itemsep=2pt, topsep=2pt]
\item ℃の0から100までは、°Fでは32から212までです。°Fの\R{目}{め}もりの\R{数}{かず}は、$212-32=180$ です。℃の100\R{目}{め}もりが、°Fの180\R{目}{め}もりにあたるのです。
\item だから、℃の1\R{目}{め}もりは、°Fの $180\div100=1.8$ \R{目}{め}もりにあたります。\R{図}{ず}でも、℃が10ふえると°Fは18ふえています。$C$\,℃なら、0℃のところから°Fの\R{目}{め}もりが $1.8\times C$ だけ\R{進}{すす}みます。
\item °Fの\R{目}{め}もりは、0℃のところで、すでに32です。だから、\R{最後}{さい|ご}に32をたします。
\end{enumerate}
これで、℃の\R{温度}{おん|ど} $C$ から°Fの\R{温度}{おん|ど} $F$ を\R{求}{もと}める\R{式}{しき}、$F=1.8C+32$ ができました。\R{確}{たし}かめると、$C=100$ なら $1.8\times100+32=212$ で、ふっとうする\R{温度}{おん|ど}と\R{合}{あ}います。

この\R{式}{しき}に $C=-10$ を\R{代入}{だい|にゅう}します。マイナスの\R{数}{かず}なので、かっこをつけます。\R{計算}{けい|さん}の\R{順番}{じゅん|ばん}は「×が＋より\R{先}{さき}」です。
\bkiinline{$F=1.8\times(-10)+32=-18+32=14$}
\qty{-10}{\degreeCelsius}は14°Fです。\R{図}{ず}の\R{目}{め}もりの\R{左}{ひだり}はしとも\R{合}{あ}っています。

\bksubsection{4}{\R{式}{しき}を、かっこごと\R{入}{い}れる}

\leavevmode\label{bk:2-13:shiki}\R{文字}{も|じ}には、\R{数}{かず}だけでなく、\R{式}{しき}も\R{代入}{だい|にゅう}できます。たとえば、「1\R{秒}{びょう}より、ほんの\R{少}{すこ}しあと」の\R{落}{お}ちた\R{距離}{きょ|り}を\R{考}{かんが}えます。その「\R{少}{すこ}し」を $h$ \R{秒}{びょう}とすると、\R{時間}{じ|かん}は $1+h$ \R{秒}{びょう}です。$h$ には、0.1や0.01など、いろいろな\R{数}{かず}が\R{入}{はい}ります。$x=5t^2$ の $t$ に $1+h$ を\R{代入}{だい|にゅう}すれば、その\R{距離}{きょ|り}が\R{表}{あらわ}せます。

\noindent\begin{minipage}[c]{0.5\linewidth}\centering
\begin{tikzpicture}[x=1mm, y=1mm]
  \node[font=\large] at (0,0) {$x=5\times$};
  \draw[bkTheme, line width=1.3pt, fill=bkTheme!8] (8,-4.5) rectangle (24,4.5);
  \node[font=\large] at (25.8,3.2) {$^2$};
  \node[draw=bkBlue, line width=0.9pt, fill=white, rounded corners=1mm, font=\normalsize, minimum height=7mm] (c) at (16,16) {$(1+h)$};
  \draw[bk link] (c) -- (16,4.5);
  \node[font=\sffamily\footnotesize, anchor=west] at (27,16) {かっこごと\R{入}{い}れる};
\end{tikzpicture}
\end{minipage}\hfill
\begin{minipage}[c]{0.47\linewidth}
\R{文字}{も|じ}は「\R{数}{かず}を\R{入}{い}れる\R{箱}{はこ}」でした。$1+h$ は、2つの\R{数}{かず}をたしたひとまとまりの\R{時間}{じ|かん}なので、かっこをつけて、まるごと\R{箱}{はこ}に\R{入}{い}れます。
\end{minipage}
\bkiinline{$x=5(1+h)^2$}
\bkiinline{\fbox{かっこごと\R{入}{い}れるのが\R{代入}{だい|にゅう}です。}}

\textbf{\R{正}{ただ}しいことを、\R{数}{かず}で\R{確}{たし}かめる}\quad $h=0.1$ として、2つのやり\R{方}{かた}で\R{計算}{けい|さん}してみます。
\begin{bkcenter}
\begin{bktable}{colspec={Q[l,m,wd=34mm] Q[l,m,wd=84mm]}, rowsep=2pt, cells={font=\small}}
\SetCell{bg=bkGreen!12} かっこごと\R{入}{い}れた\R{式}{しき} & $5(1+0.1)^2=5\times1.1^2=5\times1.21=6.05$ \\
\SetCell{bg=bkGreen!12} はじめから $t=1.1$ & $5\times1.1^2=5\times1.21=6.05$\quad \R{同}{おな}じ！ \\
\SetCell{bg=bkRed!10} かっこをつけない & $5\times1+0.1^2=5+0.01=5.01$\quad ちがう！ \\
\end{bktable}
\bktabcaption{$h=0.1$ のとき。かっこをつけないと、ちがう\R{値}{あたい}になる。}
\end{bkcenter}
かっこをつけずに $5\times1+h^2$ とすると、2\R{乗}{じょう}が $h$ だけにかかり、5も1だけにかかってしまいます。\R{別}{べつ}の\R{式}{しき}になってしまうのです。\R{第}{だい}5\R{巻}{かん}では、この $5(1+h)^2$ を\R{使}{つか}って、1\R{秒後}{びょう|ご}のボールの\R{速}{はや}さを\R{計算}{けい|さん}します。

\bksubsection{5}{\R{代入}{だい|にゅう}をくり\R{返}{かえ}して、はさんで\R{調}{しら}べる}

\leavevmode\label{bk:2-13:hasamu}ものさしを\qty{15}{cm}の\R{所}{ところ}でつかんだ\R{人}{ひと}の\R{反応時間}{はん|のう|じ|かん}は、\R{何秒}{なん|びょう}でしょう。\R{表}{ひょう}には\qty{15}{cm}がありません。そこで、\R{時間}{じ|かん}を\R{少}{すこ}しずつ\R{変}{か}えて\R{代入}{だい|にゅう}し、\qty{15}{cm}（\qty{0.15}{m}）をはさむ2つの\R{時間}{じ|かん}を\R{探}{さが}します（2-8で、$\sqrt{2}$ をはさんだのと\R{同}{おな}じ\R{考}{かんが}え\R{方}{かた}です）。

\begin{bkcenter}
\begin{bktable}{colspec={Q[l,m,wd=30mm] *{3}{Q[c,m,wd=22mm]}}, rowsep=1.5pt, cells={font=\small}}
\R{時間}{じ|かん} $t$（\R{秒}{びょう}） & 0.16 & 0.17 & 0.18 \\
$5t^2$（m） & 0.128 & 0.1445 & 0.162 \\
\qty{0.15}{m}とくらべて & \R{小}{ちい}さい & \SetCell{bg=bkYellow!40}\R{小}{ちい}さい & \SetCell{bg=bkYellow!40}\R{大}{おお}きい \\
\end{bktable}
\bktabcaption{0.17\R{秒}{びょう}と0.18\R{秒}{びょう}のあいだで、\qty{0.15}{m}をこえる。}
\end{bkcenter}
0.17\R{秒}{びょう}ではまだ\qty{0.15}{m}にとどかず、0.18\R{秒}{びょう}ではこえています。だから、\R{反応時間}{はん|のう|じ|かん}は0.17\R{秒}{びょう}と0.18\R{秒}{びょう}のあいだです。もっとくわしく\R{知}{し}りたければ、0.171、0.172、…と\R{代入}{だい|にゅう}を\R{続}{つづ}けます。\R{代入}{だい|にゅう}だけでも、ほしい\R{量}{りょう}にいくらでも\R{近}{ちか}づけるのです。ただ、\R{何回}{なん|かい}も\R{計算}{けい|さん}するのは\R{大変}{たい|へん}ですね。1\R{回}{かい}の\R{計算}{けい|さん}で\R{直接}{ちょく|せつ}\R{求}{もと}める\R{方法}{ほう|ほう}は、2-14で\R{学}{まな}びます。

\begin{taisetsu}
\R{代入}{だい|にゅう}：\R{文字}{も|じ}に\R{数}{かず}を\R{入}{い}れること。はぶいていた×をもどして\R{計算}{けい|さん}する。\par\vspace{2pt}
\R{計算}{けい|さん}の\R{順番}{じゅん|ばん}は、かっこの\R{中}{なか}→\R{累乗}{るい|じょう}→×÷→＋−。\par\vspace{2pt}
\R{単位}{たん|い}ごと\R{代入}{だい|にゅう}する。そろっていないときは、\qty{60}{s/min}などをかけて（わって）そろえてから。\par\vspace{2pt}
マイナスの\R{数}{かず}・\R{分数}{ぶん|すう}は、かっこをつけて\R{代入}{だい|にゅう}する。\par\vspace{2pt}
\R{式}{しき}も、かっこごと\R{入}{い}れる。かっこごと\R{入}{い}れるのが\R{代入}{だい|にゅう}。
\end{taisetsu}

\begin{kotoba}
\small
\bkword{\R{代入}{だい|にゅう}}{\R{式}{しき}の\R{文字}{も|じ}に、\R{数}{かず}や\R{式}{しき}を\R{入}{い}れること。}
\bkword{\R{式}{しき}の\R{値}{あたい}}{\R{代入}{だい|にゅう}して\R{計算}{けい|さん}した\R{答}{こた}え。$x=3$ のとき、$2x+1$ の\R{値}{あたい}は7。}
\bkword{\R{華氏}{か|し}（°F）}{アメリカなどで\R{使}{つか}う\R{温度}{おん|ど}の\R{表}{あらわ}し\R{方}{かた}。$F=1.8C+32$}
\end{kotoba}

\bkkigocards{\bkkigocard{°F}{かしど}{\R{温度}{おん|ど}の\R{単位}{たん|い}。\R{水}{みず}がこおる\R{温度}{おん|ど}が32°F、ふっとうする\R{温度}{おん|ど}が212°F。}{degree Fahrenheit}{2-13}}
{}

%% ---------------- 例題 ----------------
\bksection{reidai}{\R{例題}{れい|だい}でたしかめよう}

\begin{reidai}[\R{代入}{だい|にゅう}して、\R{式}{しき}の\R{値}{あたい}を\R{求}{もと}める]
$x=4$ のときの\R{式}{しき}の\R{値}{あたい}は？\quad（1）\nobreak $3x+5$\quad（2）\nobreak $x^2-2x$\quad（3）\nobreak $\dfrac{12}{x}$\par
\kangaekata はぶいていた×をもどして、$x$ に4を\R{入}{い}れます。\R{計算}{けい|さん}の\R{順番}{じゅん|ばん}は、\R{累乗}{るい|じょう}→×÷→＋−です。\par
\shiki（1）$3\times4+5=12+5=17$\quad ×を\R{先}{さき}に\par
\hspace*{2.6\zw}（2）$4^2-2\times4=16-8=8$\quad 2\R{乗}{じょう}と×を\R{先}{さき}に\par
\hspace*{2.6\zw}（3）$\dfrac{12}{4}=12\div4=3$\quad \R{分数}{ぶん|すう}はわり\R{算}{ざん}（2-4）\par
\kotae（1）\nobreak 17\quad（2）\nobreak 8\quad（3）\nobreak 3
\end{reidai}

\begin{reidai}[\R{時間}{じ|かん}の\R{単位}{たん|い}をそろえて\R{代入}{だい|にゅう}する]
\R{自転車}{じ|てん|しゃ}が\qty{6}{m/s}で3\R{分}{ぷん}\R{走}{はし}りました。$x=vt$ を\R{使}{つか}って、\R{進}{すす}んだ\R{道}{みち}のりを\R{求}{もと}めましょう。\par
\kangaekata \R{速}{はや}さは「1\R{秒}{びょう}あたり」なので、\R{時間}{じ|かん}も\R{秒}{びょう}にそろえます。3\R{分}{ぷん}に \qty{60}{s/min} をかけると、minが\R{約分}{やく|ぶん}されてsになります。\par
\shiki $\qty{3}{min}\times\qty{60}{s/min}=\qty{180}{s}$\par
\hspace*{2.6\zw}$x=\qty{6}{m/s}\times\qty{180}{s}=\qty{1080}{m}$\quad sが\R{約分}{やく|ぶん}されてm\par
\kotae \qty{1080}{m}
\end{reidai}

\begingroup\tcbset{equal height group=bkiinjA}%
\noindent\begin{minipage}[t]{0.49\linewidth}
\begin{jibun}
$a=5$ のときの\R{式}{しき}の\R{値}{あたい}は？\quad（1）\nobreak $4a-3$\quad（2）\nobreak $a^2+a$\quad（3）\nobreak $\dfrac{30}{a}$
\bkjans{（1）\nobreak 17\quad（2）\nobreak 30\quad（3）\nobreak 6}
\end{jibun}
\end{minipage}\hfill
\begin{minipage}[t]{0.49\linewidth}
\begin{jibun}
\R{電車}{でん|しゃ}が\qty{20}{m/s}で2\R{分}{ふん}\R{走}{はし}りました。$x=vt$ を\R{使}{つか}って、\R{進}{すす}んだ\R{道}{みち}のりを\R{求}{もと}めましょう。
\bkjans{$\qty{2}{min}\times\qty{60}{s/min}=\qty{120}{s}$。$\qty{20}{m/s}\times\qty{120}{s}=\qty{2400}{m}$}
\end{jibun}
\end{minipage}\par
\endgroup

\begin{reidai}[\R{長}{なが}さの\R{単位}{たん|い}をそろえて\R{代入}{だい|にゅう}する]
\R{長方形}{ちょう|ほう|けい}の\R{花}{か}だんの\R{面積}{めん|せき}は $S=ab$（たて×\R{横}{よこ}）です。たてが $a=\qty{80}{cm}$、\R{横}{よこ}が $b=\qty{2}{m}$ のとき、\R{面積}{めん|せき}を\unit{cm^2}で\R{求}{もと}めましょう。\par
\kangaekata \R{面積}{めん|せき}をcm$^2$で\R{求}{もと}めるので、\R{長}{なが}さをcmにそろえます。2 mに \qty{100}{cm/m} をかけます。\par
\shiki $\qty{2}{m}\times\qty{100}{cm/m}=\qty{200}{cm}$\quad mが\R{約分}{やく|ぶん}されてcm\par
\hspace*{2.6\zw}$S=\qty{80}{cm}\times\qty{200}{cm}=\qty{16000}{cm^2}$\quad $\mathrm{cm}\times\mathrm{cm}=\mathrm{cm^2}$\par
\kotae \qty{16000}{cm^2}
\end{reidai}

\begin{reidai}[マイナスの\R{数}{かず}・\R{分数}{ぶん|すう}を\R{代入}{だい|にゅう}する]
\R{次}{つぎ}の\R{式}{しき}の\R{値}{あたい}を\R{求}{もと}めましょう。\par
\mbox{（1）$x=-2$ のとき $3x^2$}\quad\mbox{（2）$x=\dfrac{1}{3}$ のとき $6x+1$}\par
\mbox{（3）$a=-4$ のとき $5-a$ と $-a^2$}\par
\mbox{（4）$F=1.8C+32$ で、$C=-5$ のときの $F$}\par
\kangaekata マイナスの\R{数}{かず}・\R{分数}{ぶん|すう}には、かっこをつけて\R{代入}{だい|にゅう}します。$-a^2$ は「$a^2$ にマイナスをつけたもの」です。\par
\shiki（1）$3\times(-2)^2=3\times4=12$\par
\hspace*{2.6\zw}（2）$6\times\dfrac{1}{3}+1=2+1=3$\par
\hspace*{2.6\zw}（3）$5-(-4)=5+4=9$\quad（マイナスをひく＝プラスをたす。2-2）\par
\hspace*{4.6\zw}$-(-4)^2=-\{(-4)\times(-4)\}=-16$\quad（2\R{乗}{じょう}が\R{先}{さき}）\par
\hspace*{2.6\zw}（4）$F=1.8\times(-5)+32=-9+32=23$\par
\kotae（1）\nobreak 12\quad（2）\nobreak 3\quad（3）\nobreak 9、$-16$\quad（4）\nobreak 23
\end{reidai}

\begingroup\tcbset{equal height group=bkiinjB}%
\noindent\begin{minipage}[t]{0.49\linewidth}
\begin{jibun}
\R{長方形}{ちょう|ほう|けい}の\R{板}{いた}の、たてが\qty{30}{cm}、\R{横}{よこ}が\qty{1.2}{m}です。$S=ab$ に\R{単位}{たん|い}をそろえて\R{代入}{だい|にゅう}し、\R{面積}{めん|せき}を\unit{cm^2}で\R{求}{もと}めましょう。
\bkjans{$\qty{1.2}{m}\times\qty{100}{cm/m}=\qty{120}{cm}$。$\qty{30}{cm}\times\qty{120}{cm}=\qty{3600}{cm^2}$}
\end{jibun}
\end{minipage}\hfill
\begin{minipage}[t]{0.49\linewidth}
\begin{jibun}
\R{次}{つぎ}の\R{式}{しき}の\R{値}{あたい}を\R{求}{もと}めましょう。\par
（1）\nobreak $x=-3$ のとき $x^2+2x$\par
（2）\nobreak $x=\dfrac{1}{2}$ のとき $4x^2$\par
（3）\nobreak $a=-3$ のとき $2-a$ と $-a^2$\par
（4）\nobreak $F=1.8C+32$ で、$C=-15$ のときの $F$
\bkjans{（1）\nobreak $9-6=3$\quad（2）\nobreak 1\quad（3）\nobreak 5、$-9$\quad（4）\nobreak 5}
\end{jibun}
\end{minipage}\par
\endgroup

\begin{reidai}[\R{式}{しき}をかっこごと\R{代入}{だい|にゅう}する]
（1）$y=3x+2$ の $x$ に $a+1$ を\R{代入}{だい|にゅう}して、\R{式}{しき}を\R{簡単}{かん|たん}にしましょう。\par
（2）$x=5t^2$ の $t$ に $2+h$ を\R{代入}{だい|にゅう}した\R{式}{しき}を\R{書}{か}き、$h=0.1$ のときの\R{値}{あたい}を\R{求}{もと}めましょう。\par
\kangaekata \R{式}{しき}はひとまとまりにして、かっこごと\R{代入}{だい|にゅう}します。\par
\shiki（1）$y=3(a+1)+2$\quad $x$ の\R{所}{ところ}に $(a+1)$ を\R{書}{か}く\par
\hspace*{4.6\zw}$=3a+3+2$\quad \R{分配法則}{ぶん|ぱい|ほう|そく}でかっこをはずす（2-12）\par
\hspace*{4.6\zw}$=3a+5$\quad \R{数}{かず}どうしをまとめる\par
\hspace*{2.6\zw}（2）$x=5(2+h)^2$\quad $t$ の\R{所}{ところ}に $(2+h)$ を\R{書}{か}く\par
\hspace*{4.6\zw}$h=0.1$ のとき、$5\times(2+0.1)^2=5\times2.1^2=5\times4.41=22.05$\par
\kotae（1）\nobreak $y=3a+5$\quad（2）\nobreak $x=5(2+h)^2$、22.05
\end{reidai}

\begin{reidai}[\R{代入}{だい|にゅう}をくり\R{返}{かえ}して、はさんで\R{調}{しら}べる]
アメリカの\R{天気予報}{てん|き|よ|ほう}で「\R{気温}{き|おん}は80°F」と\R{言}{い}っていました。$F=1.8C+32$ の $C$ に26と27を\R{代入}{だい|にゅう}して、80°Fが\R{何}{なん}℃と\R{何}{なん}℃のあいだか\R{調}{しら}べましょう。\par
\kangaekata \R{代入}{だい|にゅう}して\R{出}{で}た°Fを、80とくらべます。80より\R{小}{ちい}さい\R{値}{あたい}と\R{大}{おお}きい\R{値}{あたい}が\R{出}{で}たら、そのあいだです。\par
\shiki $C=26$ のとき\quad $F=1.8\times26+32=46.8+32=78.8$\par
\hspace*{9\zw}→\enspace 80より\R{小}{ちい}さい\par
\hspace*{2.6\zw}$C=27$ のとき\quad $F=1.8\times27+32=48.6+32=80.6$\par
\hspace*{9\zw}→\enspace 80より\R{大}{おお}きい\par
\kotae \qty{26}{\degreeCelsius}と\qty{27}{\degreeCelsius}のあいだ
\end{reidai}

\begingroup\tcbset{equal height group=bkiinjC}%
\noindent\begin{minipage}[t]{0.49\linewidth}
\begin{jibun}
（1）$y=2x-1$ の $x$ に $b+3$ を\R{代入}{だい|にゅう}して、\R{式}{しき}を\R{簡単}{かん|たん}にしましょう。\par
（2）$x=5t^2$ の $t$ に $3+h$ を\R{代入}{だい|にゅう}した\R{式}{しき}を\R{書}{か}き、$h=0.1$ のときの\R{値}{あたい}を\R{求}{もと}めましょう。
\bkjans{（1）\nobreak $2(b+3)-1=2b+5$\quad（2）\nobreak $x=5(3+h)^2$、$5\times3.1^2=48.05$}
\end{jibun}
\end{minipage}\hfill
\begin{minipage}[t]{0.49\linewidth}
\begin{jibun}
「\R{気温}{き|おん}は60°F」でした。$F=1.8C+32$ の $C$ に15と16を\R{代入}{だい|にゅう}して、60°Fが\R{何}{なん}℃と\R{何}{なん}℃のあいだか\R{調}{しら}べましょう。
\bkjans{$C=15$ で59°F、$C=16$ で60.8°F。\qty{15}{\degreeCelsius}と\qty{16}{\degreeCelsius}のあいだ}
\end{jibun}
\end{minipage}\par
\endgroup

\begin{reidai}[いろいろな\R{公式}{こう|しき}に、\R{単位}{たん|い}ごと\R{代入}{だい|にゅう}する]
（1）\R{円周}{えん|しゅう}＝\R{直径}{ちょっ|けい}×\R{円周率}{えん|しゅう|りつ}（2-9）を、\R{円周}{えん|しゅう}を $\ell$（エル）、\R{直径}{ちょっ|けい}を $d$ として $\ell=\pi d$ と\R{書}{か}きます。\R{直径}{ちょっ|けい}\qty{50}{cm}のタイヤの\R{円周}{えん|しゅう}を\R{求}{もと}めましょう（$\pi=3.14$）。\par
（2）\R{正方形}{せい|ほう|けい}の\R{面積}{めん|せき}は $S=a^2$（1\R{辺}{ぺん}×1\R{辺}{ぺん}）です。1\R{辺}{ぺん}が\qty{1.5}{m}の\R{正方形}{せい|ほう|けい}の\R{面積}{めん|せき}を\R{求}{もと}めましょう。\par
\kangaekata $\pi$ はただの\R{数}{かず}（3.14）で、\R{単位}{たん|い}はありません。$a^2$ は $a\times a$ なので、\R{単位}{たん|い}も2\R{回}{かい}かけます。\par
\shiki（1）$\ell=3.14\times\qty{50}{cm}=\qty{157}{cm}$\quad \R{単位}{たん|い}はcmのまま\par
\hspace*{2.6\zw}（2）$S=(\qty{1.5}{m})^2=\qty{1.5}{m}\times\qty{1.5}{m}=\qty{2.25}{m^2}$\quad $\mathrm{m}\times\mathrm{m}=\mathrm{m^2}$\par
\kotae（1）\nobreak\qty{157}{cm}\quad（2）\nobreak\qty{2.25}{m^2}
\end{reidai}

\begin{jibun}
（1）\R{直径}{ちょっ|けい}\qty{20}{cm}のお\R{皿}{さら}のふちの\R{長}{なが}さ（\R{円周}{えん|しゅう}）を、$\ell=\pi d$ で\R{求}{もと}めましょう（$\pi=3.14$）。（2）1\R{辺}{ぺん}が\qty{0.4}{m}の\R{正方形}{せい|ほう|けい}のタイルの\R{面積}{めん|せき}を、$S=a^2$ で\R{求}{もと}めましょう。
\bkjans{（1）\nobreak $3.14\times\qty{20}{cm}=\qty{62.8}{cm}$\quad（2）\nobreak $(\qty{0.4}{m})^2=\qty{0.16}{m^2}$}
\end{jibun}

%% ---------------- まちがえやすいところ ----------------
\bksection{tsumazuki}{まちがえやすいところをチェック}

\begin{tsumazuki}{×をもどしわすれる}
\bkng{$v=5$、$t=4$ のとき、$vt=54$}
\bkok{$v=5$、$t=4$ のとき、$vt=5\times4=20$}
$vt$ は $v\times t$ のこと。\R{数}{かず}を\R{入}{い}れたら、×をもどします。54は「\R{五十四}{ご|じゅう|よん}」になってしまいます。
\end{tsumazuki}

\begin{tsumazuki}{2\R{乗}{じょう}より\R{先}{さき}に、かけ\R{算}{ざん}をする}
\bkng{$t=3$ のとき、$5t^2=(5\times3)^2=15^2=225$}
\bkok{$t=3$ のとき、$5t^2=5\times3^2=5\times9=45$}
2\R{乗}{じょう}がかかっているのは $t$ だけ。2\R{乗}{じょう}を\R{先}{さき}に\R{計算}{けい|さん}します。
\end{tsumazuki}

\begin{tsumazuki}{マイナスの\R{数}{かず}に、かっこをつけない}
\bkng{$x=-3$ のとき、$x^2=-3^2=-9$}
\bkok{$x=-3$ のとき、$x^2=(-3)^2=(-3)\times(-3)=9$}
$-3^2$ は $-(3\times3)$ のこと。$-3$ を2\R{回}{かい}かけるには、かっこがいります。
\end{tsumazuki}

\begin{tsumazuki}{$-a^2$ と $(-a)^2$ をまちがえる}
\bkng{$a=2$ のとき、$-a^2=(-2)^2=4$}
\bkok{$a=2$ のとき、$-a^2=-(2\times2)=-4$}
$-a^2$ の2\R{乗}{じょう}は $a$ だけにかかります。$a$ に2を\R{入}{い}れて $-(2^2)$ です。
\end{tsumazuki}

\begin{tsumazuki}{\R{単位}{たん|い}をそろえずに\R{代入}{だい|にゅう}する}
\bkng{\qty{5}{m/s}で2\R{分}{ふん}\R{進}{すす}むと、$5\times2=\qty{10}{m}$}
\bkok{$\qty{2}{min}\times\qty{60}{s/min}=\qty{120}{s}$ にして、$\qty{5}{m/s}\times\qty{120}{s}=\qty{600}{m}$}
\R{単位}{たん|い}ごと\R{書}{か}けば、m/sとminでsが\R{消}{き}えないことに\R{気}{き}づけます。
\end{tsumazuki}

\begin{tsumazuki}{\R{式}{しき}に、かっこをつけない}
\bkng{$5t^2$ の $t$ に $1+h$ を\R{代入}{だい|にゅう}して、$5\times1+h^2$}
\bkok{$5t^2$ の $t$ に $1+h$ を\R{代入}{だい|にゅう}して、$5(1+h)^2$}
かっこがないと、$1+h$ の\R{一部}{いち|ぶ}にしか2\R{乗}{じょう}がかかりません。$h=0.1$ で\R{確}{たし}かめると、5.01と6.05でちがいます。
\end{tsumazuki}

%% ---------------- 物理・工学 ----------------
\begin{butsuri}
\leavevmode\label{bk:2-13:butsuri}\noindent\begin{minipage}[c]{0.64\linewidth}
\par\addvspace{3pt}\textbf{\R{車}{くるま}の「あと\R{何}{なん}km\R{走}{はし}れる」}\par\noindent \R{車}{くるま}のメーターには、「あと\R{何}{なん}km\R{走}{はし}れるか」が\R{表示}{ひょう|じ}されます。\R{車}{くるま}のコンピュータは、タンクに\R{残}{のこ}っているガソリンの\R{量}{りょう} $a$（L）と、\qty{1}{L}で\R{走}{はし}れる\R{距離}{きょ|り} $f$（km/L）をはかり、
\bkiinline{$d=af$}
に\R{代入}{だい|にゅう}して、\R{走}{はし}れる\R{距離}{きょ|り} $d$（\textbf{d}istance）を\R{計算}{けい|さん}しています。\R{単位}{たん|い}ごと\R{代入}{だい|にゅう}すると、
\bkiinline{$d=\qty{20}{L}\times\qty{15}{km/L}=\qty{300}{km}$}
Lが\R{約分}{やく|ぶん}されて\R{消}{き}え、kmが\R{残}{のこ}ります。
\end{minipage}\hfill
\begin{minipage}[c]{0.33\linewidth}\centering
\begin{tikzpicture}[x=1mm, y=1mm]
  \fill[bkInk!85, rounded corners=3mm] (-2,-4) rectangle (40,26);
  \bkdisplay{(4,8)}{30}{300 km}
  \node[font=\sffamily\footnotesize, text=white, anchor=south] at (19,19.5) {あと\R{走}{はし}れる\R{距離}{きょ|り}};
  \fill[bkOrange] (6,1) rectangle ++(14,2.4);
  \draw[white, line width=0.6pt] (6,1) rectangle (34,3.4);
  \node[font=\sffamily\fontsize{7}{8}\selectfont, text=white] at (19,-1.6) {ガソリンの\R{量}{りょう}};
\end{tikzpicture}
\end{minipage}

\vspace{2pt}
ガソリンがへると $a$ が\R{変}{か}わり、\R{走}{はし}り\R{方}{かた}で $f$ も\R{変}{か}わります。コンピュータは、\R{新}{あたら}しい\R{数}{かず}を\R{何度}{なん|ど}も\R{代入}{だい|にゅう}しなおして、\R{表示}{ひょう|じ}を\R{変}{か}えています。\R{天気予報}{てん|き|よ|ほう}のコンピュータも、\R{空気}{くう|き}の\R{温度}{おん|ど}や\R{風}{かぜ}の\R{速}{はや}さを、たくさんの\R{式}{しき}に\R{次々}{つぎ|つぎ}に\R{代入}{だい|にゅう}して、\R{明日}{あした}の\R{天気}{てん|き}を\R{計算}{けい|さん}しています。\R{式}{しき}は、\R{数}{かず}を\R{入}{い}れかえれば\R{何度}{なん|ど}でも\R{使}{つか}える「\R{計算}{けい|さん}の\R{道具}{どう|ぐ}」なのです。
\end{butsuri}

\begin{column}{°Fの\R{目}{め}もりは、どう\R{決}{き}めた？}
300\R{年}{ねん}ほど\R{前}{まえ}、\R{科学者}{か|がく|しゃ}のファーレンハイトは、\R{氷}{こおり}と\R{塩}{しお}をまぜて\R{作}{つく}った、とても\R{冷}{つめ}たい\R{温度}{おん|ど}を0とし、\R{人}{ひと}の\R{体温}{たい|おん}をおよそ96とする\R{目}{め}もりを\R{作}{つく}りました。これが°F（\R{華氏}{か|し}）のもとです。のちに、\R{水}{みず}がこおる\R{温度}{おん|ど}を32°F、ふっとうする\R{温度}{おん|ど}を212°Fと\R{決}{き}め\R{直}{なお}しました。\R{体温}{たい|おん}に\R{近}{ちか}い\qty{37}{\degreeCelsius}を $F=1.8C+32$ に\R{代入}{だい|にゅう}すると、$1.8\times37+32=98.6$ で、およそ99°Fです。\R{日本}{に|ほん}の℃は、スウェーデンの\R{科学者}{か|がく|しゃ}セルシウスが\R{考}{かんが}えた\R{目}{め}もりがもとで、\R{水}{みず}がこおる\R{温度}{おん|ど}とふっとうする\R{温度}{おん|ど}のあいだを100に\R{分}{わ}けています。
\end{column}

%% ---------------- 章末問題 ----------------
\bkendproblems
\bkproblemgroup{\R{基本}{き|ほん}}{1}

\bkproblem $x=3$ のときの、\R{式}{しき}の\R{値}{あたい}を\R{求}{もと}めましょう。\par
（1）\nobreak $4x-7$\quad（2）\nobreak $x^2+1$\quad（3）\nobreak $\dfrac{18}{x}$\quad（4）\nobreak $2x^2$
\bkpans{（1）\nobreak 5\quad（2）\nobreak 10\quad（3）\nobreak 6\quad（4）\nobreak 18}

\bkproblem \R{式}{しき}の\R{値}{あたい}を\R{求}{もと}めましょう。\par
（1）\nobreak $a=-2$ のとき\enspace ア\enspace $5a$\quad イ\enspace $a^2$\quad ウ\enspace $3-a$\quad エ\enspace $-a^2$\par
（2）\nobreak $a=\dfrac{1}{4}$ のとき\enspace ア\enspace $8a$\quad イ\enspace $16a^2$
\bkpans{（1）\nobreak ア $-10$\enspace イ 4\enspace ウ 5\enspace エ $-4$\quad（2）\nobreak ア 2\enspace イ 1}

\bkproblem \R{歩}{ある}く\R{人}{ひと}が\qty{1.5}{m/s}で4\R{分}{ぷん}\R{歩}{ある}きました。$x=vt$ に\R{単位}{たん|い}ごと\R{代入}{だい|にゅう}して、\R{歩}{ある}いた\R{道}{みち}のりを\R{求}{もと}めましょう。
\bkpans{$\qty{4}{min}\times\qty{60}{s/min}=\qty{240}{s}$。$\qty{1.5}{m/s}\times\qty{240}{s}=\qty{360}{m}$}

\bkproblem $F=1.8C+32$ を\R{使}{つか}って、\R{次}{つぎ}の\R{温度}{おん|ど}を°Fで\R{表}{あらわ}しましょう。\par
（1）\nobreak \qty{25}{\degreeCelsius}\quad（2）\nobreak \qty{-20}{\degreeCelsius}
\bkpans{（1）\nobreak 77°F\quad（2）\nobreak $-4$°F}

\bkproblem $y=3x-4$ の $x$ に $t-2$ を\R{代入}{だい|にゅう}して、\R{式}{しき}を\R{簡単}{かん|たん}にしましょう。
\bkpans{$3(t-2)-4=3t-6-4=3t-10$}

\bkproblem \R{長方形}{ちょう|ほう|けい}の\R{面積}{めん|せき}は $S=ab$ です。たてが $a=\qty{4}{cm}$、\R{横}{よこ}が $b=\qty{0.5}{m}$ のとき、\R{単位}{たん|い}をそろえてから\R{代入}{だい|にゅう}して、\R{面積}{めん|せき}を\unit{cm^2}で\R{求}{もと}めましょう。
\bkpans{$\qty{0.5}{m}\times\qty{100}{cm/m}=\qty{50}{cm}$。$\qty{4}{cm}\times\qty{50}{cm}=\qty{200}{cm^2}$}

\bkproblemgroup{チャレンジ}{2}

\bkproblem ものさしキャッチ（$x=5t^2$）で、\R{反応時間}{はん|のう|じ|かん}が0.15\R{秒}{びょう}の\R{人}{ひと}は、ものさしの\R{何}{なん}cmの\R{所}{ところ}でつかみますか。
\bkpans{$5\times0.15^2=\qty{0.1125}{m}$。$\qty{0.1125}{m}\times\qty{100}{cm/m}=\qty{11.25}{cm}$}

\bkproblem $x=5t^2$ の $t$ に $1+h$ を\R{代入}{だい|にゅう}した\R{式}{しき}で、$h=0.01$ のときの\R{値}{あたい}を\R{求}{もと}めましょう。また、$h=0.1$ のとき（6.05）とくらべて、5に\R{近}{ちか}づいたか\R{遠}{とお}ざかったかを\R{答}{こた}えましょう。
\bkpans{$5\times1.01^2=5.1005$。5に\R{近}{ちか}づいた}

\bkproblem \noindent\begin{minipage}[t]{0.6\linewidth}
アメリカの\R{天気予報}{てん|き|よ|ほう}で「\R{今日}{きょう}の\R{最高気温}{さい|こう|き|おん}は100°F」と\R{言}{い}っていました。$F=1.8C+32$ の $C$ に37と38を\R{代入}{だい|にゅう}して、100°Fが\R{何}{なん}℃と\R{何}{なん}℃のあいだか\R{調}{しら}べましょう。\R{右}{みぎ}の\R{温度計}{おん|ど|けい}の、\R{左}{ひだり}は℃、\R{右}{みぎ}は°Fの\R{目}{め}もりです。
\end{minipage}\hfill
\begin{minipage}[t]{0.32\linewidth}\vspace{0mm}\centering
\begin{tikzpicture}[x=1mm, y=1mm]
  \filldraw[fill=white, draw=bkGray, line width=0.8pt, rounded corners=2.4mm] (-2.4,-2) rectangle (2.4,29);
  \fill[bkRed] (0,-4) circle (3.6); \fill[bkRed] (-1,-2) rectangle (1,14.9);
  \foreach \c in {30,35,40,45} {\pgfmathsetmacro\yy{(\c-30)*1.4+4}\draw[bkInk, line width=0.6pt] (-2.4,\yy) -- (-4.4,\yy); \node[font=\sffamily\footnotesize, anchor=east] at (-5,\yy) {\c};}
  \foreach \f in {90,100,110} {\pgfmathsetmacro\yy{((\f-32)/1.8-30)*1.4+4}\draw[bkInk, line width=0.6pt] (2.4,\yy) -- (4.4,\yy); \node[font=\sffamily\footnotesize, anchor=west] at (5,\yy) {\f};}
  \node[font=\sffamily\footnotesize, anchor=south] at (-7,30) {℃};
  \node[font=\sffamily\footnotesize, anchor=south] at (7,30) {°F};
\end{tikzpicture}
\end{minipage}
\bkpans{$C=37$ で98.6°F、$C=38$ で100.4°F。\qty{37}{\degreeCelsius}と\qty{38}{\degreeCelsius}のあいだ}

\begin{furikaeri}
\begin{checklist}
\item \R{公式}{こう|しき}の\R{文字}{も|じ}に\R{数}{かず}を\R{代入}{だい|にゅう}して、\R{式}{しき}の\R{値}{あたい}を\R{求}{もと}められる。\hfill{\footnotesize\sffamily\color{bkTheme}→\pageref{bk:2-13:dainyu}ページ}
\item \R{単位}{たん|い}ごと\R{代入}{だい|にゅう}し、\R{単位}{たん|い}をそろえて\R{計算}{けい|さん}できる。\hfill{\footnotesize\sffamily\color{bkTheme}→\pageref{bk:2-13:tani}ページ}
\item マイナスの\R{数}{かず}・\R{分数}{ぶん|すう}・\R{式}{しき}を、かっこごと\R{代入}{だい|にゅう}できる。\hfill{\footnotesize\sffamily\color{bkTheme}→\pageref{bk:2-13:kakko}ページ}
\end{checklist}
\end{furikaeri}

\bkfukidashi{delta}{100°Fが\R{何}{なん}℃かを、37と38を\R{代入}{だい|にゅう}して、はさんで\R{調}{しら}べたね。\R{式}{しき}を\R{変形}{へん|けい}すれば、ほしい\R{量}{りょう}を\R{直接}{ちょく|せつ}とり\R{出}{だ}せるよ。\R{次}{つぎ}の2-14で\R{学}{まな}ぼう。}
